% Generated draft: requires mathematical and source review.
\documentclass[11pt]{article}
\usepackage[margin=25mm]{geometry}
\usepackage{fontspec}
\setmainfont{Latin Modern Roman}
\usepackage{amsmath,amssymb,mathtools,unicode-math}
\setmathfont{Latin Modern Math}
\usepackage{enumitem,needspace,xurl,hyperref}
\setlength{\parindent}{0pt}
\setlength{\parskip}{5pt}
\setlength{\emergencystretch}{3em}
\newcommand{\N}{\mathbb{N}}
\newcommand{\Z}{\mathbb{Z}}
\newcommand{\Q}{\mathbb{Q}}
\newcommand{\R}{\mathbb{R}}
\newcommand{\C}{\mathbb{C}}
\newcommand{\F}{\mathbb{F}}
\newcommand{\A}{\mathbb{A}}
\newcommand{\PP}{\mathbb{P}}
\newcommand{\E}{\mathbb{E}}
\newcommand{\eps}{\varepsilon}
\newcommand{\dd}{\,\mathrm{d}}
\begin{document}
\section*{UnsolvedMath ID: 801096}
\textbf{Source problem code:} catalogue entry
\subsection{Definitions and mathematical statement}
For a rational parameter \(c\in\Q\), define the polynomial map
\[
f_c\colon\Q^2\longrightarrow\Q^2,\qquad f_c(x,y)=(y,y^2+c+x).
\]
Write \(f_c^n\) for its \(n\)-fold iterate. A point \(P\in\Q^2\) is periodic if \(f_c^n(P)=P\) for some integer \(n\geq1\); its least period is
\[
\operatorname{per}_{f_c}(P)=\min\{n\geq1:f_c^n(P)=P\}.
\]
\textbf{Conjecture 48, under the arithmetic interpretation below.} For every \(c\in\Q\) and every periodic point \(P\in\Q^2\),
\[
\operatorname{per}_{f_c}(P)\in\{1,2,3,4,6,8\}.
\]
Here the source phrase ``over \(\Q\)'' is provisionally interpreted as requiring both a rational parameter and a rational affine point. The supplied excerpt does not explicitly state these quantifiers. Likewise, ``period \(N\)'' is interpreted as least period, rather than merely the condition \(f_c^N(P)=P\). These interpretations require human confirmation against the source context. The conjecture excludes other least periods; it does not assert that every listed period occurs.

The introductory source paragraph discusses dependence of an unspecified bound on \(d\) and on \([K:\Q]\), with adjoining periodic points to the base field offered as the reason for the latter dependence. Neither the bound nor \(d\), \(K\), or the preceding hypotheses are defined in the supplied excerpt, so that contextual assertion cannot be formulated precisely here.
\subsection{Short English statement}
Under the provisional interpretation that \(c\) and both point coordinates are rational, every periodic point of \((x,y)\mapsto(y,y^2+c+x)\) should have least period in \(\{1,2,3,4,6,8\}\). The source leaves the parameter and point quantifiers implicit. This is a statement-only draft requiring human mathematical review.
\subsection{Sources}
\begin{itemize}
\item Local UnsolvedMath registry, permanent ID \(801096\), source problem code \texttt{catalogue entry}, in the AMR Open Problem Lists collection. The supplied record labels the problem open and lists Berger and Bedford as proposers, with proposed year 2023; these are supplied registry labels, not independently verified assessments.
\item Berger and Bedford, \emph{Henon Maps: A List of Open Problems} (2023), Conjecture 48. Supplied source URL: \url{https://arxiv.org/abs/2312.03907}. The record reports access on 2026-10-16 and extraction by \texttt{arxiv\_latex}; these provenance details are supplied and have not been independently verified here.
\item The supplied attribution/context is ``see [I]''. No bibliographic entry identifying [I] is included, so its authorship and relationship to this conjecture remain unresolved.
\end{itemize}
\end{document}
